
\documentclass[12pt,a4paper]{article}

% Packages
\usepackage[utf8]{inputenc}
\usepackage{amsmath}
\usepackage{graphicx}
\usepackage{booktabs}
\usepackage{xcolor}
\usepackage{hyperref}
\usepackage{geometry}

% Page settings
\geometry{margin=1in}

% Document information
\title{\textbf{Introduction to LaTeX}}
\author{Vishesh}
\date{\today}

\begin{document}

\maketitle

\begin{abstract}
This report introduces LaTeX, a document preparation
system commonly used for academic and technical writing.
It demonstrates document formatting, mathematical
equations, tables, lists and references.
\end{abstract}

\tableofcontents
\newpage

\section{Introduction}

LaTeX is a document preparation system used to create
professional documents. It is especially useful for
scientific writing, research papers and technical reports.

Unlike ordinary word processors, LaTeX uses commands
to control the structure and formatting of documents.

\section{Important Features}

Some important features of LaTeX are:

\begin{itemize}
    \item Professional document formatting
    \item Mathematical equation support
    \item Automatic numbering of sections
    \item Tables and figures
    \item Bibliography and references
\end{itemize}

\section{Mathematical Equations}

LaTeX provides excellent support for mathematical
expressions.

For example, the quadratic formula is:

\begin{equation}
x = \frac{-b \pm \sqrt{b^2 - 4ac}}{2a}
\end{equation}

Another example is:

\begin{equation}
E = mc^2
\end{equation}

\section{Tables}

The following table presents some commonly used
LaTeX commands.

\begin{center}
\begin{tabular}{ll}
\toprule
\textbf{Command} & \textbf{Purpose} \\
\midrule
\texttt{\textbackslash section\{\}} & Creates a section \\
\texttt{\textbackslash textbf\{\}} & Makes text bold \\
\texttt{\textbackslash begin\{itemize\}} & Creates a list \\
\texttt{\textbackslash includegraphics} & Inserts an image \\
\bottomrule
\end{tabular}
\end{center}

\section{Conclusion}

Through this practical, we learned how to create
a structured academic report using LaTeX.

We also explored document formatting, mathematical
equations, tables and lists.

\end{document}
